\documentclass[11pt,a4paper]{article}
\usepackage[utf8]{inputenc}
\usepackage[margin=1in]{geometry}
\usepackage{graphicx}
\usepackage{booktabs}
\usepackage{amsmath}
\usepackage{hyperref}
\usepackage{float}

\title{Deep Reinforcement Learning Assignment 3\\Meta and Transfer Learning}
\author{Ben-Gurion University of the Negev}
\date{January 2024}

\begin{document}

\maketitle

\section{Introduction}

This report presents the implementation and analysis of meta and transfer learning techniques for Deep Reinforcement Learning. We implement an Actor-Critic architecture and evaluate three transfer learning approaches across three OpenAI Gymnasium environments: CartPole-v1, Acrobot-v1, and MountainCarContinuous-v0.

\textbf{Transfer Learning Methods Implemented:}
\begin{itemize}
    \item \textbf{Section 1:} Individual Actor-Critic training from scratch (baseline)
    \item \textbf{Section 2:} Fine-tuning pre-trained models with output layer re-initialization
    \item \textbf{Section 3:} Progressive Networks with frozen source networks and lateral connections
\end{itemize}

\section{Architecture Design}

\subsection{Standardized Input/Output Dimensions}

To enable transfer learning across different environments, we standardize all input and output dimensions:
\begin{itemize}
    \item \textbf{Observation dimension:} 6 (padded with zeros for smaller environments)
    \item \textbf{Action dimension:} 3 (with action masking for environments with fewer actions)
\end{itemize}

\subsection{Network Architecture}

Both Actor and Critic networks use identical architectures:
\begin{itemize}
    \item Input layer: 6 units (standardized observation)
    \item Hidden layer 1: 128 units with ReLU activation
    \item Hidden layer 2: 128 units with ReLU activation
    \item Output layer: 3 units (Actor) / 1 unit (Critic)
\end{itemize}

Weight initialization uses Xavier/Glorot uniform initialization for improved gradient flow.

\section{Section 1: Individual Training Results}

\subsection{Training Configuration}
\begin{itemize}
    \item Episodes: 1500
    \item Learning rate: 0.001
    \item Discount factor ($\gamma$): 0.99
    \item Entropy coefficient: 0.01
    \item Value loss coefficient: 0.5
\end{itemize}

\subsection{Results}

\begin{table}[H]
\centering
\caption{Section 1: Individual Training Results}
\begin{tabular}{lccccc}
\toprule
\textbf{Environment} & \textbf{Episodes} & \textbf{Time (s)} & \textbf{Final Avg Reward} & \textbf{Target} & \textbf{Converged} \\
\midrule
CartPole-v1 & 1500 & 11.60 & 21.18 & $>$475 & No \\
Acrobot-v1 & 1500 & 255.15 & -500.00 & $>$-100 & No \\
MountainCarContinuous-v0 & 1500 & 470.78 & -62.94 & $>$90 & No \\
\bottomrule
\end{tabular}
\end{table}

\textbf{Observations:} None of the environments achieved convergence with the baseline Actor-Critic implementation. This establishes a baseline for comparing transfer learning approaches.

\section{Section 2: Fine-Tuning Results}

\subsection{Fine-Tuning Procedure}
\begin{enumerate}
    \item Load the fully trained model from the source environment
    \item Re-initialize the output layer weights (fc3) of both Actor and Critic
    \item Train the network on the target environment
\end{enumerate}

This approach preserves learned feature representations in hidden layers while adapting the decision layer to the new task.

\subsection{Results}

\begin{table}[H]
\centering
\caption{Section 2: Fine-Tuning Results}
\begin{tabular}{lcccc}
\toprule
\textbf{Transfer Pair} & \textbf{Episodes} & \textbf{Time (s)} & \textbf{Final Avg Reward} & \textbf{Improvement} \\
\midrule
Acrobot $\rightarrow$ CartPole & 1500 & 11.12 & 21.87 & +0.69 \\
CartPole $\rightarrow$ MountainCar & 1500 & 466.44 & -62.04 & +0.90 \\
\bottomrule
\end{tabular}
\end{table}

\textbf{Comparison with Baseline:}
\begin{itemize}
    \item CartPole: Fine-tuning from Acrobot improved reward by +0.69 (21.18 $\rightarrow$ 21.87)
    \item MountainCar: Fine-tuning from CartPole improved reward by +0.90 (-62.94 $\rightarrow$ -62.04)
\end{itemize}

Fine-tuning showed marginal improvements but did not achieve convergence, suggesting the hidden layer representations from source tasks may not transfer optimally to these specific target tasks.

\section{Section 3: Progressive Networks Results}

\subsection{Progressive Network Architecture}

Progressive Networks use multiple frozen source networks with lateral connections to a trainable target network:

\begin{itemize}
    \item \textbf{Source networks:} Frozen (non-trainable), provide hidden layer features
    \item \textbf{Lateral connections:} Transform source hidden features (128-dim $\rightarrow$ 128-dim)
    \item \textbf{Target network:} Trainable, combines own features with lateral inputs
\end{itemize}

The lateral connections allow the target network to selectively use relevant features from source tasks while learning new task-specific representations.

\subsection{Results}

\begin{table}[H]
\centering
\caption{Section 3: Progressive Networks Results}
\begin{tabular}{lcccc}
\toprule
\textbf{Sources $\rightarrow$ Target} & \textbf{Episodes} & \textbf{Time (s)} & \textbf{Final Avg Reward} & \textbf{Improvement} \\
\midrule
\{Acrobot, MountainCar\} $\rightarrow$ CartPole & 1500 & 14.81 & 22.53 & +1.35 \\
\{CartPole, Acrobot\} $\rightarrow$ MountainCar & 1500 & 638.60 & -62.35 & +0.59 \\
\bottomrule
\end{tabular}
\end{table}

\textbf{Comparison with Other Methods:}

\begin{table}[H]
\centering
\caption{Summary: All Methods Comparison}
\begin{tabular}{lccc}
\toprule
\textbf{Target Environment} & \textbf{Baseline} & \textbf{Fine-Tuned} & \textbf{Progressive} \\
\midrule
CartPole-v1 & 21.18 & 21.87 (+0.69) & \textbf{22.53 (+1.35)} \\
MountainCarContinuous-v0 & -62.94 & \textbf{-62.04 (+0.90)} & -62.35 (+0.59) \\
\bottomrule
\end{tabular}
\end{table}

\section{Analysis and Discussion}

\subsection{Transfer Learning Effectiveness}

\textbf{Progressive Networks showed the best improvement for CartPole} (+1.35 over baseline), suggesting that combining features from multiple source tasks can provide richer representations than single-source fine-tuning.

\textbf{For MountainCar, fine-tuning performed slightly better} than progressive networks (+0.90 vs +0.59), possibly because the continuous control task benefits more from adapting a single well-trained policy than combining discrete task features.

\subsection{Convergence Challenges}

None of the approaches achieved the target convergence thresholds. Potential factors:
\begin{itemize}
    \item \textbf{Episode count:} 1500 episodes may be insufficient for convergence
    \item \textbf{Hyperparameter tuning:} Default parameters may not be optimal for all environments
    \item \textbf{MountainCar discretization:} Continuous action space discretized to 3 actions limits policy expressiveness
    \item \textbf{Reward shaping:} MountainCarContinuous has sparse rewards, making learning difficult
\end{itemize}

\subsection{Key Findings}

\begin{enumerate}
    \item Transfer learning consistently improved performance over training from scratch
    \item Progressive networks provided the largest improvement for discrete action spaces (CartPole)
    \item Fine-tuning was more effective for continuous control tasks (MountainCar)
    \item Lateral connections in progressive networks allow selective feature reuse
\end{enumerate}

\section{Conclusion}

This assignment demonstrated the implementation and evaluation of transfer learning techniques for reinforcement learning. While none of the approaches achieved full convergence, both fine-tuning and progressive networks showed consistent improvements over baseline training.

Progressive networks showed particular promise for combining knowledge from multiple source tasks, with the best overall improvement observed for CartPole-v1. The modular architecture of progressive networks, with frozen source columns and trainable lateral connections, provides a principled approach to avoiding catastrophic forgetting while enabling knowledge transfer.

Future work could explore:
\begin{itemize}
    \item Extended training with more episodes
    \item Hyperparameter optimization for each environment
    \item Different network architectures for source-target matching
    \item Continuous action spaces with actor networks outputting distribution parameters
\end{itemize}

\end{document}
